\documentclass[a4paper,12pt]{article}
\usepackage[spanish,activeacute]{babel}
\usepackage[latin1]{inputenc}

%% Para las columnas
\usepackage{multicol}


\usepackage{amssymb, amsmath}
\usepackage{latexsym}
\usepackage{multicol}


%% Para numeraci?n autom?tica
\newcounter{nroproblema}
\setcounter{nroproblema}{1}
\newcounter{nroapartado}
\setcounter{nroapartado}{1}
\newcommand{\n}{%
    \vspace{.1in}\noindent{\bf \arabic{nroproblema}.\ }%
    \addtocounter{nroproblema}{1}%
    \setcounter{nroapartado}{1}%
    }%
\newcommand{\apartado}{\vspace*{0.2cm}\hspace*{0.2cm}%
   {\bf \alph{nroapartado})}
   \addtocounter{nroapartado}{1}%
   }

%% m?genes

\setlength{\unitlength}{1mm} \addtolength{\voffset}{-15mm}
\addtolength{\textheight}{30mm} \addtolength{\hoffset}{-15mm}
\addtolength{\textwidth}{30mm}


\begin{document}


\begin{flushright}
1${}^o$ de grado en F�sica, curso 2013/14

\today
\end{flushright}


\begin{center}{\bf {\Large \'ALGEBRA I. HOJA 5}}\end{center}

 \vskip.35cm




Problemas de libro Linear Algebra Done Wrong (LADW) de Treil: 
\begin{itemize}
\item p.46 2.1-2.2;
\item  p.51 3.1-3.9;
\item  pp.55-56, 5.1-5.6;
\item  pp 58-59 6.1-6.2.
\end{itemize}
% ;
% pp. 66-68 7.1-7.14;
% pp. 71-72 8.1-8.6.


 
   \n Resolver los siguientes
    sistemas de ecuaciones:
    $$
    a)\left\{\begin{array}{rrr}
      x+2y-z &= & 7 \\
      2x+y+z &= & 6 \\
      x-y+3z &= & -1
    \end{array}
    \right.\hspace{5mm}b)\left\{\begin{array}{rrr}
      2x+3y-5z-2 &= & 0 \\
      3x-y+2z+1 &= & 0 \\
      5x+4y-6z-3 &= & 0
    \end{array}
    \right.\hspace{5mm}c)\left\{\begin{array}{rrr}
      2x+y+4z+8t &= & -1 \\
      x+3y-6z+2t &= & 3 \\
      3x-2y+2z-2t &= & 8 \\
      2x-y+2z &= & 4
    \end{array}
    \right.
    $$
    Se pueden comprobar los resultados obtenidos por substituci�n
    o aplicando el m�todo de Gauss.

    Sol: a)(5/3,8/3,0), b)(-1/5,14/5,6/5), c)(2,-3,-3/2,1/2).

   \n Hallar los valores de $a$ para que los siguientes sistemas sean compatibles indeterminados.



$$a) \left. \begin{array}{r@{\hspace{0pt}}r@{\hspace{0pt}}r@{\hspace{0pt}}c}
ax&+y&+z&=1\\
x&+ay&+z&=1\\
x&+y&+az&=1\end{array}\right\} \qquad
b) \left. \begin{array}{r@{\hspace{0pt}}r@{\hspace{0pt}}r@{\hspace{0pt}}c}
ax& +y&+2z&=-2\\
2x&+y&+az&=3\\
x&+ay&+2z&=-2
\end{array}\right\}
\qquad c) \left. \begin{array}{r@{\hspace{0pt}}r@{\hspace{0pt}}r@{\hspace{0pt}}c}
-x&+ay&+az&=0\\
ax&-y&+az&=0\\
ax&+ay&-z&=0
\end{array}\right\}$$
sol: a) $a=-2$ \'o $a=1$ \quad b) $a=-3$ \'o $a=1$ \quad c) $a=-1$ \'o $a= \frac{1}{2}$


\n Estudiar los rangos de las siguientes matrices como funci�n de
$\lambda$:
$$\left(\begin{matrix}
2 &3 & 1& 7\\
3&  7 & -6 &-2\\
5& 8& 1 & \lambda
\end{matrix}\right)
, \left(\begin{matrix}
3& 1& 1& 4\\
\lambda &4& 10& 1\\
1& 7& 17& 3\\
2& 2& 4& 3 \end{matrix}\right) , \left(\begin{matrix}
1 & \lambda & -1& 2\\
2& -1 &\lambda&  5\\
1 &10& -6& 1 \end{matrix}\right)$$



\n Encontrar bases de los siguientes subespacios:
$$\left.\begin{array}{c@{\hspace{3pt}}c@{\hspace{3pt}}c@{\hspace{3pt}}c@{\hspace{3pt}}c@{\hspace{3pt}}c}
x_1 &&+x_3 +&2x_4 &= 0 \\
x_1 &&-x_3&& = 0 \end{array}\right\}\subset
\mathbb{R}^4$$
$$\left.\begin{array}{c@{\hspace{3pt}}c@{\hspace{3pt}}c@{\hspace{3pt}}c@{\hspace{3pt}}c@{\hspace{3pt}}c@{\hspace{3pt}}c}
x_1 &+x_2& +x_3 &+x_4 &+x_5 &= 0\\
x_1& -x_2 &+x_3& -x_4& +x_5 &= 0\end{array}\right\}\subset
\mathbb{R}^5 \qquad
\left.\begin{array}{c@{\hspace{3pt}}c@{\hspace{3pt}}c@{\hspace{3pt}}c@{\hspace{3pt}}c@{\hspace{3pt}}c}
3x_1 &-x_2 &+3x_3& -x_4 &+3x_5 &= 0\\
3x_1& +x_2 &+3x_3 &+x_4 &+3x_5 &= 0\end{array}\right\}\subset
\mathbb{R}^5$$

\vspace{0.2cm}

\n Encontrar una base del subespacio $S
\subset \mathcal{M}_{2\times2} (\mathbb{R})$ definido por
$$S = \left\{\left( \begin{matrix}a &b \\
c & d\end{matrix} \right)\Bigg|
\begin{array}{c@{\hspace{3pt}}c@{\hspace{3pt}}c@{\hspace{3pt}}c@{\hspace{3pt}}c}
2a &+ b&- c &+ d &= 0\\
a& + b &+ c& - d &= 0\end{array}\right\}$$

\vspace{0.2cm}

\n Encontrar una base del espacio vectorial de los polinomios de grado menor o igual que tres divisibles por $x- 1$.

\vspace{0.2cm}

\n Comprobar que el subespacio de $\mathbb{R}^4$ generado por
$\{(1, 0, 1, -1)^T, (1, -1, 1, -1)^T\}$ coincide con el espacio de
soluciones del sistema:
$$
\left.\begin{array}{c@{\hspace{3pt}}c@{\hspace{3pt}}c@{\hspace{3pt}}c@{\hspace{3pt}}c}
x_1 &\ &+x_3 &+2x_4 &= 0\\
x_1 &&-x_3 &&= 0\end{array}\right\}\subset
\mathbb{R}^4.$$

% \vspace{0.2cm}


\n Hallar las ecuaciones cartesianas de los siguientes subespacios de $\mathbb{R}^4$:
$$\begin{array}{c@{\hspace{2pt}}l}
S_1 =& \operatorname{Span}\{(1, 0, 1, 0)^T,(2, 1, 0, -1)^T, (1, -1, 3, 1)^T\},\\[+2pt]
S_2 = & \operatorname{Span}\{(3, 1, 0, -1)^T,(1, 1, -1, -1)^T, (7, 1, 2, -1)^T\},\\[+2pt]
S_3 = & \operatorname{Span}\{(0, 2, 5, 0)^T\}.
\end{array}$$

% \vspace{0.2cm}



\n Siendo $S_1 =
\operatorname{Span}\{(1, -5, 2, 0)^T, (1, -1, 0, 2)^T\}$ y
$S_2 = \mathcal L\{(3, -5, 2, 1)^T, (2, 0, 0, 1)^T\}$,
hallar una base
y las ecuaciones cartesianas de $S_1 + S_2$ y de $S_1\cap S_2$.

\vspace{0.2cm}

\n Sean
\begin{align*}
S_1 & =
\operatorname{Span}\{(1, 0, 1, 0, 1)^T,(2, 1, 0, -1, 0)^T,(2, 0, 1, 0, 1)^T\}, \\
% \text{y} \quad
S_2 & =
\operatorname{Span}\{(3, 1, 0, -1, 0)^T,(1, 1, -1, -1, -1)\}^T.
\end{align*}
Comprobar que $S_1 + S_2 = S_1$ y $ S_1\cap S_2 = S_2$.

\vspace{0.2cm}

\n Siendo $S_1 =
\operatorname{Span}\{(0, 1, 1, 0)^T, (1, 0, 0, 1)^T\}$ y
$$S_2 \equiv\left\{\begin{array}{c@{\hspace{3pt}}c@{\hspace{3pt}}c@{\hspace{3pt}}c@{\hspace{3pt}}c@{\hspace{3pt}}c}
x_1& +x_2& +x_3& +x_4 &= 0 \\
x_1 &+x_2 &&&= 0\end{array} \right.$$
hallar una base y las ecuaciones cartesianas de $S_1 + S_2$ y $S_1\cap S_2$.

\vspace{0.2cm}

\n   Comprobar que son complementarios los subespacios definidos por las ecuaciones:
$$S_1 \equiv
\left\{\begin{array}{c@{\hspace{3pt}}c@{\hspace{3pt}}c@{\hspace{3pt}}c@{\hspace{3pt}}c@{\hspace{3pt}}c}
-2x_1&-5x_2& +4x_3 &&= 0\\
-2x_1 &+7x_2 &&+4x_4 &= 0 \end{array}\right.
\qquad
S_2 \equiv \left\{\begin{array}{c@{\hspace{3pt}}c@{\hspace{3pt}}c@{\hspace{3pt}}c@{\hspace{3pt}}c@{\hspace{3pt}}c}
x_1 &+2x_2&& +x_4& = 0\\
x_1 & -2x_2 &+x_3&& = 0
\end{array}\right.$$

\vspace{0.2cm}

\n Hallar una base y las ecuaciones cartesianas de un espacio complementario de

\apartado $S_1 =
\operatorname{Span}\{(0, 2, 5, 0)^T,(-1, 1, 3, 2)^T\}$.

\apartado $S_2 =
\operatorname{Span}\{(1, 1, 0, 0)^T,(1, 0, 1, 0)^T,(0, 0, 1, 1)^T,(0, 1, 0, 1)^T\}$.
\vspace{0.2cm}

\n Hallar una base de $S_1 \cap S_2$ y de $S_1 + S_2$, donde
$$S_1 =
\operatorname{Span} \left\{
\left( \begin{matrix}
1 &1 \\
0 &0
\end{matrix}\right),
\left( \begin{matrix}
0 &0 \\
1 &1
\end{matrix}\right)\right\}\, ,  \qquad
S_2 =
\operatorname{Span} \left\{
\left( \begin{matrix}
1 &0 \\
0 &1
\end{matrix}\right),
\left( \begin{matrix}
0 &1\\
1 &0
\end{matrix}\right)\right\}.$$
Hallar tambi\'en las ecuaciones cartesianas de $S_1\cap S_2$ y de $S_1 + S_2$.

\vspace{0.2cm}

%%%%%%%%%%%%%%%%%%%%%%%%%%%%%%%%%%%%%%%%%%%%%%%

          %\newpage 

%%%%%%%%%%%%%%%%%%%%%%%%%%%%%%%%%%%%%%%%%%%%%%%

\n En $\mathcal{M}_{2\times 2} (\mathbb{R})$ se consideran los subespacios
que conmutan con cada una de las siguientes
matrices:
$$A =\left(\begin{matrix}
1& 1\\
0 &2\end{matrix}\right)\, ;
\qquad B =\left(\begin{matrix}
1 &1\\
1 &1\end{matrix}\right).$$

\apartado Hallar una base de cada uno de ellos, de su suma y de su intersecci\'on.

\apartado Hallar una base de un espacio complementario del subespacio de las matrices que conmutan con $A$.

\bigskip

\n Hallar los subespacios suma e intersecci\'on de los siguientes subespacios de las matrices
cuadradas de orden $n$:

\apartado El subespacio de las matrices triangulares superiores de orden $n$ y el subespacio de las matrices
triangulares inferiores.

\apartado El subespacio de las matrices sim\'etricas de orden $n$ y el subespacio de las matrices antisim\'etricas de orden $n$.

\bigskip

\n Consideramos los subespacios del espacio de polinomios de grado $\leq 3$ $S_1$, formado por los polinomios m\'ultiplos de $x+1$, y $S_2$, formado por los polinomios m\'ultiplos de $x-1$. Hallar los subespacios suma e intersecci\'on de $S_1$ y $S_2$.

\bigskip

\n En $\mathbb{R}^4$ sean $U =
\operatorname{Span}\{u_1 , u_2 \}$ y $V = \operatorname{Span}\{v_1 , v_2 \}$ donde
$$
u_1 = (1, 1, 2,-\lambda)^T, \quad u_2 = (-1, 1, 0, -\lambda)^T, \quad
v_1 = (1, \lambda, 2, -\lambda)^T, \quad v_2 = (2, 3, \lambda, 1)^T.
$$
Hallar seg\'un los valores de $\lambda$ las dimensiones de $U$ , $V$,
$U + V$ y  $U \cap  V$.



\end{document}
